% Complete Section 4.5 and Box 4.1, original PDF 54-60 / printed 44-50.
% The printed signs and the last summand in Box 4.1 are retained, not corrected.
% block: hodge-r3-heading page:54 type:section
\subsection{{Example: $\mathbb R^3$ 中的 Hodge Laplacian}}\label{sec:hodge-r3}

% block: hodge-r3-introduction pages:54,55 type:ordinary-paragraph
{现在回到我们贯穿使用的例子, 即 $\mathbb R^3$ 中具有 Lipschitz 边界的区域上的 de Rham 复形. 在 Section~\ref{sec:l2-de-rham-complex} 中已看到, 它构成闭 Hilbert 复形. 因此, 对每个 $k=0,1,2,3$, 都得到与该复形相伴的 Hodge Laplace 问题, 它是一个适定的偏微分方程边值问题. 为了解释这些问题, 只需解释空间 $V^k$ 和 $V_k^*$ 以及算子 $d^k$ 和 $d_k^*$; 为此我们将频繁使用定义域复形 \eqref{eq:l2-de-rham-domain-complex} 及对偶定义域复形 \eqref{eq:l2-dual-domain-complex}.}

% block: hodge-r3-k0-heading page:55 type:subsection
\subsubsection{$k=0$}\label{subsec:hodge-r3-k0}

% block: hodge-r3-k0-domain page:55 type:ordinary-paragraph
{假设区域连通, 则调和形式空间 $\mathsf H^0$ 仅由常数构成, 且 $P_{\mathsf H}f=\bar f$, 即 $f$ 在 $\Omega$ 上的平均值. 由于 $d_0^*=0$, 有 $L^0u=(dd^*+d^*d)u=d^*du=-\operatorname{div}\operatorname{grad}u=-\Delta u$, 定义域为}
\[
\begin{aligned}
D(L)&=\{\,u\in H^1(\Omega)\mid\operatorname{grad}u\in\mathring H(\operatorname{div},\Omega)\,\}\\
&=\{\,u\in H^1(\Omega)\mid\Delta u\in L^2(\Omega),\ \partial u/\partial n=0\text{ on }\partial\Omega\,\},
\end{aligned}
\]
{其中 $\Delta$ 表示通常的 Laplacian $\Delta=\sum_{i=1}^3\partial^2/\partial x_i^2$. 给定 $f\in L^2(\Omega)$, 强表述 \eqref{eq:hodge-strong-formulation} 就是求 $u\in H^1$, 使得}
\begin{equation}\label{eq:hodge-r3-neumann-strong}
-\Delta u=f-\bar f\text{ in }\Omega,\qquad\frac{\partial u}{\partial n}=0\text{ on }\partial\Omega,\qquad\int u=0.
\end{equation}

% block: hodge-r3-neumann-interpretation page:55 type:ordinary-paragraph
{因此, 与 de Rham 复形相伴的 $k=0$ Hodge Laplace 问题, 就是 Poisson 方程通常的 Neumann 问题; 减去 $\bar f$ 的调整是为了使数据 $f$ 与齐次 Neumann 边界条件相容.}

% block: hodge-r3-neumann-weak page:55 type:ordinary-paragraph
{接下来讨论原始和混合弱表述. 由于 $V^{-1}=0$, 当 $k=0$ 时, 混合弱表述 \eqref{eq:hodge-mixed-formulation} 中的 $\sigma$ 和 $\tau$ 均消失. 因此在 $k=0$ 时两种弱表述没有区别. 它们都是求 $u\in H^1$, 使得}
\begin{equation}\label{eq:hodge-r3-neumann-weak}
\int\operatorname{grad}u\cdot\operatorname{grad}v=\int(f-\bar f)v,\qquad v\in H^1,\qquad\int_\Omega u=0.
\end{equation}

% block: hodge-r3-neumann-equivalence page:55 type:ordinary-paragraph
{由 Theorem~\ref{thm:hodge-formulation-equivalence} 已知, 此问题等价于强表述问题 \eqref{eq:hodge-r3-neumann-strong}. 为使讨论具体, 现在直接验证 \eqref{eq:hodge-r3-neumann-weak} 蕴含 \eqref{eq:hodge-r3-neumann-strong}, 另一个方向更容易. 在 \eqref{eq:hodge-r3-neumann-weak} 中取 $v\in C_0^\infty(\Omega)$, 可见}
\begin{equation}\label{eq:hodge-r3-neumann-distribution}
-\operatorname{div}\operatorname{grad}u=f-\bar f
\end{equation}
{在分布意义下成立, 因此由 $f\in L^2$ 可知, $\operatorname{grad}u$ 属于 $H(\operatorname{div})$. 故可在分部积分公式 \eqref{eq:normal-trace-green} 中取 $w=\operatorname{grad}u$, 再利用 \eqref{eq:hodge-r3-neumann-weak} 和 \eqref{eq:hodge-r3-neumann-distribution}, 得到}
\[
\langle\gamma_n(\operatorname{grad}u),\operatorname{tr}v\rangle=\int(\operatorname{div}\operatorname{grad}u)v+\int\operatorname{grad}u\cdot\operatorname{grad}v=0,\quad v\in H^1.
\]
{由于 $\operatorname{tr}v$ 可以是 $H^{1/2}(\partial\Omega)$ 的任意元素, 可知 $\gamma_n(\operatorname{grad}u)=0$, 即边界上有 $\partial u/\partial n=0$. 注意, 在弱表述中, 边界条件是\emph{自然}施加的, 即通过分部积分由方程蕴含, 而不是施加在求解 $u$ 所在的空间中.}

% block: hodge-r3-k0-boundary-problems page:55 type:ordinary-paragraph
{再看 $\mathcal B$ 和 $\mathcal B^*$ 问题, 由于 $\mathcal B^0=0$, 并没有新的内容. 因此, 在 $k=0$ 时不存在 $\mathcal B$ 问题, 而 $\mathcal B^*$ 问题本质上就是完整的 Hodge Laplacian 问题.}

% block: hodge-r3-k0-summary page:56 type:ordinary-paragraph
{概括 $k=0$ 的情形, $0$-形式的 Hodge Laplacian 就是 Poisson 方程的 Neumann 问题. 此时原始和混合弱表述一致, 都是 Neumann 问题的标准弱表述, 其中边界条件是\emph{自然}施加的.}

% block: hodge-r3-k1-heading page:56 type:subsection
\subsubsection{$k=1$}\label{subsec:hodge-r3-k1}

% block: hodge-r3-k1-strong page:56 type:ordinary-paragraph
{$k=1$ 的情形更有意思. 调和形式空间 $\mathsf H^1$ 是向量场空间, 其维数为 $b_1$, 即第一个 Betti 数. 强表述是求 $u\in H(\operatorname{curl})\cap\mathring H(\operatorname{div})$, 满足 $\operatorname{curl}u\in\mathring H(\operatorname{curl})$, $\operatorname{div}u\in H^1$, 以及偏微分方程}
\begin{equation}\label{eq:hodge-r3-vector-laplace}
\operatorname{curl}\operatorname{curl}u-\operatorname{grad}\operatorname{div}u=f-P_{\mathsf H}f\text{ in }\Omega.
\end{equation}
{注意, 微分算子 $\operatorname{curl}\operatorname{curl}-\operatorname{grad}\operatorname{div}$ 是负的向量 Laplacian $-\vec\Delta$, 因为对光滑向量场有 \[\operatorname{curl}\operatorname{curl}u-\operatorname{grad}\operatorname{div}u=(-\Delta u_1,-\Delta u_2,-\Delta u_3)\] 其中 $u=(u_1,u_2,u_3)$.}

% block: hodge-r3-magnetic-boundary page:56 type:ordinary-paragraph
{定义域 $D(L^1)$ 的定义蕴含边界条件}
\begin{equation}\label{eq:hodge-r3-magnetic-boundary}
u\cdot n=0,\qquad(\operatorname{curl}u)\times n=0\text{ on }\partial\Omega,
\end{equation}
{当然, 还需要附加条件 $u\perp\mathsf H^1$ 来保证唯一性. 条件 \eqref{eq:hodge-r3-magnetic-boundary} 有时称为\emph{磁边界条件}, 其原因将在 Section~\ref{sec:maxwell-equations} 讨论 Maxwell 方程时变得更清楚.}

% block: hodge-r3-k1-primal page:56 type:ordinary-paragraph
{原始弱表述是求 $u\in H(\operatorname{curl})\cap\mathring H(\operatorname{div})$, 使得}
\begin{equation}\label{eq:hodge-r3-magnetic-primal}
\begin{gathered}
\langle\operatorname{curl}u,\operatorname{curl}v\rangle+\langle\operatorname{div}u,\operatorname{div}v\rangle=\langle f-P_{\mathsf H}f,v\rangle,\\
v\in H(\operatorname{curl})\cap\mathring H(\operatorname{div}).
\end{gathered}
\end{equation}
{注意, 在此表述中, 边界条件 $u\cdot n=0$ 是本质的 (施加在求解 $u$ 所在的空间中), 而条件 $\operatorname{curl}u\times n=0$ 是自然的.}

% block: hodge-r3-k1-mixed page:56 type:ordinary-paragraph
{相比之下, 混合弱表述是求 $\sigma\in H^1$, $u\in H(\operatorname{curl})$, $p\in\mathsf H^1$, 不施加任何边界条件, 并满足}
\begin{equation}\label{eq:hodge-r3-magnetic-mixed}
\begin{aligned}
\int\sigma\tau-\int u\cdot\operatorname{grad}\tau&=0,&\tau&\in H^1,\\
\int\operatorname{grad}\sigma\cdot v+\int\operatorname{curl}u\cdot\operatorname{curl}v+\int p\cdot v&=\int f\cdot v,&v&\in H(\operatorname{curl}),\\
\int u\cdot q&=0,&q&\in\mathsf H^1.
\end{aligned}
\end{equation}
{在此表述中, 边界条件 \eqref{eq:hodge-r3-magnetic-boundary} 都是自然的, 分别来自第一个和第二个方程中的分部积分. 注意, 除未知向量场 $u$ 外, 此表述还以标量函数 $\sigma=-\operatorname{div}u$ 为第二个未知量 (当然还有 $p=P_{\mathsf H}f$).}

% block: hodge-r3-k1-b-problem page:56 type:ordinary-paragraph
{当 $f\in\mathcal B^1$, 即 $f$ 为梯度时, 解 $u$ 也是梯度, 这可以表述为 $\operatorname{curl}u=0$ 且 $u\perp\mathsf H^1$. 因此 $\mathcal B^1$ 问题是求 $u$, 使得}
\[
-\operatorname{grad}\operatorname{div}u=f,\ \operatorname{curl}u=0\text{ in }\Omega,\qquad u\cdot n=0\text{ on }\partial\Omega,\qquad u\perp\mathsf H^1.
\]

% block: hodge-r3-div-curl-application page:57 type:ordinary-paragraph
{由抽象理论可知, 对 $f\in\mathcal B^1$, 此问题具有唯一解; 而且, 可通过对同一 $f$ 求解完整的 Hodge Laplacian 问题 (例如通过混合弱表述) 来计算此解. 此类问题的一个应用是求解 div-curl 问题}
\[
\operatorname{div}u=F,\qquad\operatorname{curl}u=0,
\]
{这可以通过取 $f=\operatorname{grad}F$ 实现 (为与 $u$ 的边界条件相容, 需要 $\int F=0$).}

% block: hodge-r3-k1-bstar-problem page:57 type:ordinary-paragraph
{另一方面, 当 $f\in\mathcal B_1^*$, 即 $f\in\operatorname{curl}\mathring H(\operatorname{curl})$ 时, 就得到 $\mathcal B^*$ 问题. 此时 $u$ 由下式确定:}
\[
\begin{gathered}
\operatorname{curl}\operatorname{curl}u=f,\ \operatorname{div}u=0\text{ in }\Omega,\\
u\cdot n=0,\ \operatorname{curl}u\times n=0\text{ on }\partial\Omega,\qquad u\perp\mathsf H^1.
\end{gathered}
\]
{这种 curl-curl 问题出现在电磁学中. 同样, 可通过求解完整的 Hodge Laplace 问题来确定其解.}

% block: hodge-r3-k2-heading page:57 type:subsection
\subsubsection{$k=2$}\label{subsec:hodge-r3-k2}

% block: hodge-r3-k2-domain page:57 type:ordinary-paragraph
{此情形仍得到向量 Laplace 方程 \eqref{eq:hodge-r3-vector-laplace}, 但边界条件不同 (调和形式也不同, 其中 $\dim\mathsf H^2=b_2$). $L^2$ 的定义域由满足 $\operatorname{div}u\in\mathring H^1$ 和 $\operatorname{curl}u\in H(\operatorname{curl})$ 的 $u\in H(\operatorname{div})\cap\mathring H(\operatorname{curl})$ 构成. 因此得到\emph{电边界条件}:}
\[
u\times n=0,\ \operatorname{div}u=0\text{ on }\partial\Omega.
\]

% block: hodge-r3-k2-mixed page:57 type:ordinary-paragraph
{这些条件中, 第一个在原始变分表述中是本质的, 第二个是自然的. 当然, 在混合弱表述中两者都是自然的; 该表述是求 $\sigma\in H(\operatorname{curl})$, $u\in H(\operatorname{div})$, $p\in\mathsf H^2$, 使得}
\begin{equation}\label{eq:hodge-r3-electric-mixed}
\begin{aligned}
\int\sigma\cdot\tau-\int u\cdot\operatorname{curl}\tau&=0,&\tau&\in H(\operatorname{curl}),\\
\int\operatorname{curl}\sigma\cdot v+\int\operatorname{div}u\operatorname{div}v+\int p\cdot v&=\int f\cdot v,&v&\in H(\operatorname{div}),\\
\int u\cdot q&=0,&q&\in\mathsf H^2.
\end{aligned}
\end{equation}
{注意, 此表述与 $k=1$ 时的混合弱表述 \eqref{eq:hodge-r3-magnetic-mixed} 很不相同. 这里的附加场 $\sigma$ 是向量场, 而在 $k=1$ 时它是标量场.}

% block: hodge-r3-k2-boundary-problems page:57 type:ordinary-paragraph
{当 $k=2$ 时, $\mathcal B$ 问题为}
\[
\operatorname{curl}\operatorname{curl}u=f,\ \operatorname{div}u=0\text{ in }\Omega,\qquad u\times n=0,\text{ on }\partial\Omega,\qquad u\perp\mathsf H^2,
\]
{若 $f\in\operatorname{curl}H(\operatorname{curl})$, 则它具有唯一解. $\mathcal B_2^*$ 问题为}
\[
\begin{gathered}
-\operatorname{grad}\operatorname{div}u=f,\ \operatorname{curl}u=0\text{ in }\Omega,\\
u\times n=0,\ \operatorname{div}u=0\text{ on }\partial\Omega,\qquad u\perp\mathsf H^2,
\end{gathered}
\]
{对 $f\in\operatorname{grad}\mathring H^1$ 可解. 注意, $\mathcal B^2$ 问题的微分方程与 $\mathcal B_1^*$ 问题的微分方程一致, $\mathcal B_2^*$ 问题的微分方程与 $\mathcal B^1$ 问题的微分方程一致, 但相应的边界条件并不相同.}

% block: hodge-r3-table page:58 type:table
\begin{table}[htbp]\centering\small
\caption{{de Rham 复形的 Hodge Laplace 问题. $\sigma$ 是混合弱表述中的附加变量. 带下划线的边界条件 (不涉及导数的那些条件) 在原始弱表述中是本质的, 其余是自然的. 所有边界条件在混合弱表述中都是自然的.}}\label{tab:hodge-r3-boundary-conditions}
\setlength{\tabcolsep}{4pt}
\begin{tabular}{cccll}\hline
$k$&$L^k$&$\sigma$&{边界条件}&{调和形式}\\\hline
$0$&$-\Delta$&$0$&Neumann: $\partial u/\partial n=0$&$\mathsf H^0$ $(\dim b_0)$\\\noalign{\vskip4pt}
$1$&$-\vec\Delta$&$-\operatorname{div}u$&{磁}: $\underline{u\cdot n=0}$, $\operatorname{curl}u\times n=0$&$\mathsf H^1$ $(\dim b_1)$\\\noalign{\vskip4pt}
$2$&$-\vec\Delta$&$\operatorname{curl}u$&{电}: $\underline{u\times n=0}$, $\operatorname{div}u=0$&$\mathsf H^2$ $(\dim b_2)$\\\noalign{\vskip4pt}
$3$&$-\Delta$&$-\operatorname{grad}u$&Dirichlet: $\underline{u=0}$&$\mathsf H^3=0$\\\hline
\end{tabular}
\end{table}

% block: hodge-r3-k3-heading page:58 type:subsection
\subsubsection{$k=3$}\label{subsec:hodge-r3-k3}

% block: hodge-r3-k3-dirichlet page:58 type:ordinary-paragraph
{空间 $\mathcal B_3^*$ 和 $\mathsf H^3$ 都为零. Laplacian 的定义域为 $D(L^3)=\{\,u\in\mathring H^1\mid\Delta u\in L^2\,\}$, 因此该边值问题就是 Poisson 方程的 Dirichlet 问题:}
\[
-\Delta u=f\text{ in }\Omega,\qquad u=0\text{ on }\partial\Omega.
\]
{其原始弱表述是显然的: 求 $u\in\mathring H^1$, 使得}
\[
\int\operatorname{grad}u\cdot\operatorname{grad}v=\int fv,\qquad v\in\mathring H^1.
\]
{混合弱表述是求 $\sigma\in H(\operatorname{div})$, $u\in L^2$, 使得}
\[
\begin{gathered}
\int\sigma\cdot\tau-\int u\operatorname{div}\tau=0,\qquad\tau\in H(\operatorname{div}),\\
\int\operatorname{div}\sigma\,v=\int fv,\qquad v\in L^2.
\end{gathered}
\]
{在混合表述中, Dirichlet 条件是自然的.}

% block: hodge-r3-table-summary page:58 type:ordinary-paragraph
{Table~\ref{tab:hodge-r3-boundary-conditions} 汇集了 de Rham 复形的 Hodge Laplacian 问题的主要特征.}

% block: hodge-r3-boundary-duality-heading page:58 type:subsection
\subsubsection{{带边界条件的 de Rham 复形与对偶性}}\label{subsec:de-rham-boundary-duality}

% block: hodge-r3-boundary-complex page:58 type:ordinary-paragraph
{刚才讨论了抽象理论在 $L^2$ de Rham 复形中的应用, 其定义域复形由 \eqref{eq:l2-de-rham-domain-complex} 给出. 也可以将理论应用于\emph{带边界条件的 $L^2$ de Rham 复形}, 即基空间仍为 $L^2$, 而定义域复形为下式的 Hilbert 复形:}
\[
0\to\mathring H^1\xrightarrow{\operatorname{grad}}\mathring H(\operatorname{curl})\xrightarrow{\operatorname{curl}}\mathring H(\operatorname{div})\xrightarrow{\operatorname{div}}L^2\to0.
\]
{除若干符号和指标的变化外, 此 Hilbert 复形本质上就是 \eqref{eq:l2-de-rham-domain-complex} 的对偶复形, 即 \eqref{eq:l2-dual-domain-complex}. 具体地, 带边界条件的 de Rham 复形中指标为 $k$ 的空间, 就是原 de Rham 复形的对偶复形中指标为 $3-k$ 的空间. 因此, 带边界条件的 de Rham 复形的对偶复形为}
\[
0\leftarrow L^2\xleftarrow{-\operatorname{div}}H(\operatorname{div})\xleftarrow{\operatorname{curl}}H(\operatorname{curl})\xleftarrow{-\operatorname{grad}}H^1\leftarrow0,
\]
{即除符号和指标外, 就是原 de Rham 复形.}

% block: hodge-r3-boundary-table page:59 type:table
\begin{table}[htbp]\centering\small
\caption{{带边界条件的 de Rham 复形的 Hodge Laplace 问题. 在混合表述中, 所有边界条件都是本质的; 在原始表述中, 带下划线的边界条件是本质的. 见 Table~\ref{tab:hodge-r3-boundary-conditions}.}}\label{tab:hodge-r3-relative-boundary-conditions}
\setlength{\tabcolsep}{3pt}
\begin{tabular}{cccll}\hline
$k$&$L^k$&$\sigma$&{边界条件}&{调和形式}\\\hline
$0$&$-\Delta$&$0$&Dirichlet: $\underline{u=0}$&$\mathring{\mathsf H}^0=0$\\\noalign{\vskip4pt}
$1$&$-\vec\Delta$&$-\operatorname{div}u$&{电}: $\underline{u\times n=0}$, $\operatorname{div}u=0$&$\mathring{\mathsf H}^1=\mathsf H^2$ $(\dim b_2)$\\\noalign{\vskip4pt}
$2$&$-\vec\Delta$&$\operatorname{curl}u$&{磁}: $\underline{u\cdot n=0}$, $\operatorname{curl}u\times n=0$&$\mathring{\mathsf H}^2=\mathsf H^1$ $(\dim b_1)$\\\noalign{\vskip4pt}
$3$&$-\Delta$&$-\operatorname{grad}u$&Neumann: $\partial u/\partial n=0$&$\mathring{\mathsf H}^3=\mathsf H^0$ $(\dim b_0)$\\\hline
\end{tabular}
\end{table}

% block: hodge-r3-poincare-lefschetz page:59 type:ordinary-paragraph
{若用 $\mathring{\mathcal B}^k$, $\mathring{\mathcal Z}^k$ 和 $\mathring{\mathsf H}^k$ 表示带边界条件的复形的边界空间, 闭链空间和调和形式空间, 则有 $\mathring{\mathcal B}^k=\mathcal B^{3-k}$, 对 $\mathcal Z$ 和 $\mathsf H$ 也有类似关系. (在 $n$ 维流形上的关系 $\mathsf H^{n-k}=\mathring{\mathsf H}^k$ 称为 \emph{Poincaré–Lefschetz 对偶性}. 对无边界流形, 它成为 $\mathsf H^{n-k}=\mathsf H^k$, 即 Poincaré 对偶性.) 容易验证, 除指标变化外, 带边界条件的复形的强表述和原始弱表述与通常的 de Rham 复形的相应表述一致, 因而得到与上面相同的边值问题. 例如 $k=0$ 的问题就是带 Dirichlet 边界条件的 Poisson 方程. 然而, 对 $k>0$ 的问题, 混合弱表述带来了新的内容. 例如当 $k=1$ 时, 求 $\sigma\in\mathring H^1$, $u\in\mathring H(\operatorname{curl})$, $p\in\mathring{\mathsf H}^1=\mathsf H^2$, 使得}
\[
\begin{aligned}
\int\sigma\tau-\int u\cdot\operatorname{grad}\tau&=0,&\tau&\in\mathring H^1,\\
\int\operatorname{grad}\sigma\cdot v+\int\operatorname{curl}u\cdot\operatorname{curl}v+\int p\cdot v&=\int f\cdot v,&v&\in\mathring H(\operatorname{curl}).\\
\int u\cdot q&=0,&q&\in\mathsf H^2.
\end{aligned}
\]

% block: hodge-r3-relative-mixed-comparison page:59 type:ordinary-paragraph
{这是向量 Laplacian 在边界条件 $\operatorname{div}u=0$, $u\times n=0$ 下的混合表述. 这些边界条件与通常的 de Rham 复形中 $k=2$ 问题的边界条件相同. 但是, 那个问题的混合方法 \eqref{eq:hodge-r3-electric-mixed} 完全不同, 辅助变量 $\sigma$ 使用向量场而不是这里的标量场, 并且边界条件是自然施加的, 而不是像这里这样作为本质条件施加.}

% block: hodge-r3-relative-summary page:59 type:ordinary-paragraph
{对 $k=2$ 和 $k=3$ 也有类似情形, 细节留待读者推导. Table~\ref{tab:hodge-r3-relative-boundary-conditions} 给出了汇总.}

% block: hilbert-summary-heading page:60 type:ordinary-paragraph
\begin{tcolorbox}[enhanced,breakable,colback=KuangNei,frame hidden,boxrule=0pt,sharpish corners,before skip=10pt,after skip=10pt]
\refstepcounter{booksummary}\label{box:hilbert-complexes-summary}
\noindent\textbf{Box~\thebooksummary:} {Chapter~\ref{chap:hilbert-complexes} Hilbert 复形内容摘要.}\par

% block: hilbert-summary-definition page:60 type:ordinary-paragraph
{\emph{Hilbert 复形}是由 Hilbert 空间 $W^k$ 以及从一个空间到下一个空间的稠定闭算子 $d^k$ 构成的序列, 其中 $d^k$ 将其定义域映入 $d^{k+1}$ 的零空间. 若所有 $d^k$ 都具有闭像, 则称此 Hilbert 复形是\emph{闭的}.}

% block: hilbert-summary-domain-complex page:60 type:ordinary-paragraph
{令 $V^k=D(d^k)$, 得到定义域复形}
\[
0\to V^0\xrightarrow{d}V^1\xrightarrow{d}\cdots\xrightarrow{d}V^n\to0,
\]
{它是由 Hilbert 空间和有界线性算子构成的上链复形.}

% block: hilbert-summary-cocycles page:60 type:ordinary-paragraph
{用 $\mathcal Z^k$ 和 $\mathcal B^k$ 表示闭上链和上边界链, 对偶复形的相应空间则记为 $\mathcal Z_k^*$ 和 $\mathcal B_k^*$.}

% block: hilbert-summary-laplacian page:60 type:ordinary-paragraph
{\emph{Hodge Laplacian} 是从 $W^k\to W^k$ 的无界算子 $L^k$, 由 $d^*d+dd^*$ 给出并取自然定义域. 其核为\emph{调和形式}空间 $\mathsf H^k=\mathcal Z^k\cap\mathcal Z_k^*$. 对闭 Hilbert 复形, $\mathsf H^k$ 同构于上同调空间 $\mathcal H^k$.}

% block: hilbert-summary-decomposition page:60 type:ordinary-paragraph
{任一闭 Hilbert 复形都给出正交的 \emph{Hodge 分解}}
\[
W^k=\mathcal B^k\BookOrthSum\mathsf H^k\BookOrthSum\mathcal Z^k
\]
{并满足 \emph{Poincaré 不等式}}
\[
\|z\|_V\leq c_P\|dz\|,\qquad z\in\mathcal Z^{k\perp_V}.
\]

% block: hilbert-summary-solvability page:60 type:ordinary-paragraph
{Hodge Laplace 问题}
\[
L^ku=f-P_{\mathsf H}f,\qquad u\perp\mathsf H^k,
\]
{对任意 $f\in W^k$ 在 $D(L^k)$ 中具有唯一解. 此解也可以用原始弱表述 \eqref{eq:hodge-primal-formulation} 或混合弱表述 \eqref{eq:hodge-mixed-formulation} 来刻画.}
\end{tcolorbox}
